\documentclass[12pt,a4paper]{article}

% Encoding and fonts
\usepackage[utf8]{inputenc}
\usepackage[T1]{fontenc}
\usepackage{lmodern}

% Page layout
\usepackage[top=1in, bottom=1in, left=1in, right=1in]{geometry}

% Typography
\usepackage{microtype}
\usepackage{parskip}

% Language
\usepackage[english]{babel}

% Headers and footers
\usepackage{fancyhdr}
\pagestyle{fancy}
\fancyhf{}
\fancyhead[L]{\leftmark}
\fancyhead[R]{\thepage}
\fancyfoot[C]{\thepage}
\renewcommand{\headrulewidth}{0.4pt}
\renewcommand{\footrulewidth}{0pt}

% Section styling
\usepackage{titlesec}
\titleformat{\section}{\Large\bfseries}{\thesection}{1em}{}
\titleformat{\subsection}{\large\bfseries}{\thesubsection}{1em}{}
\titleformat{\subsubsection}{\normalsize\bfseries}{\thesubsubsection}{1em}{}
\titlespacing*{\section}{0pt}{2.5ex plus 1ex minus .2ex}{1.5ex plus .2ex}
\titlespacing*{\subsection}{0pt}{2ex plus 1ex minus .2ex}{1ex plus .2ex}
\titlespacing*{\subsubsection}{0pt}{1.5ex plus 1ex minus .2ex}{0.8ex plus .2ex}

% List formatting
\usepackage{enumitem}
\setlist[itemize]{topsep=0.5ex, itemsep=0.25ex, parsep=0pt}
\setlist[enumerate]{topsep=0.5ex, itemsep=0.25ex, parsep=0pt}

% Essential packages
\usepackage{graphicx}
\usepackage[table]{xcolor}
\usepackage{booktabs}
\usepackage{array}
\usepackage{tabularx}
\usepackage{longtable}
\usepackage{amsmath}
\usepackage{amssymb}
\usepackage[normalem]{ulem}
\rowcolors{2}{gray!10}{white}
\renewcommand{\arraystretch}{1.3}

% Cross-references
\usepackage{hyperref}
\usepackage{cleveref}

% Hyperref setup
\hypersetup{
    colorlinks=true,
    linkcolor=blue,
    filecolor=magenta,
    urlcolor=cyan,
}

% Code listings
\usepackage{listings}
\definecolor{codebg}{RGB}{248,248,248}
\definecolor{codeframe}{RGB}{200,200,200}
\definecolor{codekeyword}{RGB}{0,0,180}
\definecolor{codecomment}{RGB}{0,128,0}
\definecolor{codestring}{RGB}{163,21,21}
\lstset{
    basicstyle=\ttfamily\small,
    breaklines=true,
    breakatwhitespace=false,
    frame=single,
    framerule=0.5pt,
    rulecolor=\color{codeframe},
    backgroundcolor=\color{codebg},
    keepspaces=true,
    showstringspaces=false,
    tabsize=4,
    xleftmargin=1em,
    xrightmargin=1em,
    aboveskip=1em,
    belowskip=1em,
    keywordstyle=\color{codekeyword}\bfseries,
    commentstyle=\color{codecomment}\itshape,
    stringstyle=\color{codestring},
}

% YAML language definition for listings
\lstdefinelanguage{YAML}{
    keywords={true,false,null,y,n},
    keywordstyle=\color{blue}\bfseries,
    sensitive=false,
    comment=[l]{\#},
    morecomment=[s]{/*}{*/},
    commentstyle=\color{gray}\ttfamily,
    stringstyle=\color{red}\ttfamily,
    morestring=[b]',
    morestring=[b]",
}

% TOML language definition for listings
\lstdefinelanguage{TOML}{
    keywords={true,false},
    keywordstyle=\color{blue}\bfseries,
    sensitive=true,
    comment=[l]{\#},
    commentstyle=\color{gray}\ttfamily,
    stringstyle=\color{red}\ttfamily,
    morestring=[b]',
    morestring=[b]",
}

% Dockerfile language definition for listings
\lstdefinelanguage{Dockerfile}{
    keywords={FROM,RUN,CMD,LABEL,EXPOSE,ENV,ADD,COPY,ENTRYPOINT,VOLUME,USER,WORKDIR,ARG,ONBUILD,STOPSIGNAL,HEALTHCHECK,SHELL},
    keywordstyle=\color{blue}\bfseries,
    sensitive=false,
    comment=[l]{\#},
    commentstyle=\color{gray}\ttfamily,
    stringstyle=\color{red}\ttfamily,
    morestring=[b]',
    morestring=[b]",
}

% JSON language definition for listings
\lstdefinelanguage{JSON}{
    keywords={true,false,null},
    keywordstyle=\color{blue}\bfseries,
    sensitive=true,
    commentstyle=\color{gray}\ttfamily,
    stringstyle=\color{red}\ttfamily,
    morestring=[b]",
}

% Code listing float environment
\usepackage{float}
\newfloat{listing}{htbp}{lol}[section]
\floatname{listing}{Listing}

% Admonition boxes
\usepackage{tcolorbox}
\tcbuselibrary{skins,breakable}

% Admonition style definitions
\newtcolorbox{admonitionbox}[2][]{
    enhanced,
    breakable,
    colback=#2!5,
    colframe=#2!80!black,
    fonttitle=\bfseries,
    title=#1,
    left=2mm,
    right=2mm,
}

% Synthon tuple boxes
\newtcolorbox{synthonbox}{
    enhanced,
    colback=blue!5,
    colframe=blue!75!black,
    fontupper=\large,
    center upper,
    boxrule=1pt,
    arc=4pt,
}

\begin{document}

\title{SynthOmnicon: Programmable Matter Domain}
\date{\today}

\maketitle

\emph{This document applies the SynthOmnicon framework to programmable matter --- materials that can be commanded to change their macroscopic properties on demand. It follows the structure of PROTEINS.md: each material class gets an encoding, the algebra is applied, and the results are interpreted as predictions. The work was generated by \texttt{programmable\_matter\_tests1.py} (encodings, distances, meet, path algebra, Varma probe, cross-domain analogy, programmability lattice) and \texttt{programmable\_matter\_tests2.py} (Primitive Jacobian, tensor products, DesignPipeline monad, formal predictions P-38--P-47). All algebra results are exact outputs of the framework. All predictions are formally registered in \texttt{PRIMITIVE\_PREDICTIONS.md}.}

\begin{center}
\rule{0.5\textwidth}{0.4pt}
\end{center}

\subsection{Why Programmable Matter?}

Programmable matter --- materials that change shape, stiffness, optical properties, or chemical function in response to a stimulus --- is one of the central challenges of twenty-first century materials science. The design problem is always the same: you want a material that can be in at least two macroscopically distinguishable states, can be switched between them reliably, can maintain each state long enough to be useful, and can switch reversibly (or controllably irreversibly).

These requirements map onto the synthon grammar with unusual precision. ''Macroscopically distinguishable states'' = two different synthon tuples. ''Reliable switching'' = a path exists in the HotSwap graph. ''State maintenance'' = $F$ and $K$ are compatible with the required dwell time. ''Reversibility'' = directed distances $d(A \to B)$ and $d(B \to A)$ are both finite and accessible. The framework does not describe what programmable matter \emph{is} --- it describes what it \emph{can do}, in the same sense that a synthon encodes interaction affordances rather than molecular constitution.

The question ''what can this material be programmed to do?'' becomes: ''what paths exist in the HotSwap graph from this synthon's encoding?''

\begin{center}
\rule{0.5\textwidth}{0.4pt}
\end{center}

\subsection{Part I: The Eleven Encodings}

\subsubsection{Class 1 --- DNA Nanotechnology}

DNA nanotechnology provides two logically distinct programmable modes: the static addressable scaffold (DNA origami) and the dynamic reconfiguration mechanism (strand displacement).

\textbf{DNA origami (folded):} The final origami structure is locked by hundreds of staple strands simultaneously. Watson-Crick specificity gives $F_\hbar$ (each staple--scaffold pair has $K_a \gg 10^7$ M$^{-1}$). Folding requires thermal annealing over hours: $K_\text{slow}$. The complete addressable scaffold is a mesoscale object (10--100 nm): $G_\gimel$. All staple strands must be present simultaneously for the target shape to emerge: $\Gamma_\wedge(\text{SPECIFIC})$.

\[\langle D_{\triangle\wedge} \;;\; T_\text{network} \;;\; R_{\supseteq} \;;\; P_{\text{pm\_sym}} \;;\; F_\hbar \;;\; K_\text{slow} \;;\; G_\gimel \;;\; \Gamma_\wedge(\text{SPEC}) \;;\; \Phi_\text{sub} \;;\; n:m \rangle\]

\textbf{DNA strand displacement:} Toehold-mediated branch migration is a directed ratchet operating at the molecular scale. The toehold hybridizes first (sequential grammar, $\Gamma_\to$), then branch migration follows --- each event is local and operates one displacement at a time ($G_\beth$). Toehold $K_a \sim 10^4$--$10^6$ M$^{-1}$ places this at $F_\eth$. Branch migration barrier \textasciitilde{}60--80 kJ/mol: $K_\text{mod}$.

\[\langle D_\wedge \;;\; T_\text{linear} \;;\; R_{\supseteq} \;;\; P_{\text{pm\_sym}} \;;\; F_\eth \;;\; K_\text{mod} \;;\; G_\beth \;;\; \Gamma_\to(\text{SEL}) \;;\; \Phi_\text{sub} \;;\; 1:1 \rangle\]

\textbf{Distance between the two states:} $d = 6.10$ (largest programmability pair distance of the five classes). Meet conflicts on $D$, $T$, and $\Gamma$ --- the mechanism changes completely between the two modes, not just the thermodynamic parameters. This encodes a deep design insight: DNA origami and strand displacement are not two modes of the same system, they are two different primitive architectures that happen to share the same material (DNA). Combining them in a single device requires constructing the tensor product --- a higher-level system that contains both grammars as coupled subsystems.

\subsubsection{Class 2 --- Colloidal Assembly}

\textbf{Colloidal crystal ($T < T_m$):} Long-range crystalline order locked by cooperative van der Waals, depletion, and electrostatic forces. At the ensemble level, collective binding gives $F_\hbar$ (many contacts per particle). Nucleation barrier requires slow annealing: $K_\text{slow}$. Global crystalline order propagates without spatial attenuation: $G_\aleph$. Any particle from the batch can join any lattice site: $\Gamma_\wedge(\text{BROAD})$.

\[\langle D_\triangle \;;\; T_{\text{network\_sym}} \;;\; R_{\supseteq} \;;\; P_{\text{pm\_sym}} \;;\; F_\hbar \;;\; K_\text{slow} \;;\; G_\aleph \;;\; \Gamma_\wedge(\text{BRD}) \;;\; \Phi_\text{sub} \;;\; n:n \rangle\]

\textbf{Colloidal fluid ($T > T_m$):} Thermal fluctuations overcome all local barriers: $K_\text{fast}$. Individual contacts are weak: $F_\ell$. Only nearest-neighbour correlations remain: $G_\beth$. Any nearest neighbour suffices: $\Gamma_\vee(\text{BROAD})$.

\[\langle D_\triangle \;;\; T_\text{network} \;;\; R_{\supseteq} \;;\; P_{\text{pm\_sym}} \;;\; F_\ell \;;\; K_\text{fast} \;;\; G_\beth \;;\; \Gamma_\vee(\text{BRD}) \;;\; \Phi_\text{sub} \;;\; n:n \rangle\]

\textbf{Distance:} $d = 5.10$. Meet conflicts on $T$ and $\Gamma$: the topology changes from the centrosymmetric crystal lattice ($T_{\text{network\_sym}}$) to a generic disordered network ($T_\text{network}$). This $T$-conflict encodes a first-order melting transition --- colloidal melting has latent heat and coexistence, consistent with the framework's prediction that $T$-conflicting pairs are first-order (P-40).

\subsubsection{Class 3 --- Biomolecular Condensates (LLPS)}

\textbf{Condensate liquid (dynamic droplet):} IDR-IDR interactions via $\pi$-$\pi$, cation-$\pi$, and electrostatic contacts. Each individual contact is weak ($F_\ell$: individually \textasciitilde{}1--3 $k_BT$). Highly dynamic: FRAP recovery $t_{1/2} \sim 10$--30 s ($K_\text{fast}$). The mesoscale droplet (100 nm--10 $\mu$m) defines the scale of structural coherence ($G_\gimel$). Promiscuous partner set: $\Gamma_\vee(\text{BROAD})$. The condensate is encoded as $\Phi_c$ --- near the spinodal/critical point. This is the defining claim: condensates that are near their critical concentration do not simply phase-separate, they operate in the G/D degenerate regime where molecular-scale inputs propagate to the droplet scale.

\[\langle D_{\triangle\wedge} \;;\; T_\text{network} \;;\; R_{\supseteq} \;;\; P_{\text{pm\_sym}} \;;\; F_\ell \;;\; K_\text{fast} \;;\; G_\gimel \;;\; \Gamma_\vee(\text{BRD}) \;;\; \Phi_c \;;\; n:m \rangle\]

\textbf{Condensate gel (arrested/solid state):} The pathological transition from dynamic droplet to solid gel is encoded by three simultaneous primitive changes: $F_\ell \to F_\hbar$ (contacts lock --- no FRAP recovery), $K_\text{fast} \to K_\text{trap}$ (kinetically arrested; multiple misfolding pathways), $G_\gimel \to G_\aleph$ (system-spanning amyloid-like network --- the arrest is global). $\Phi$ drops from $\Phi_c$ to $\Phi_\text{sub}$: the system moves away from criticality into the arrested state. $\Gamma$ changes from $\Gamma_\vee$ (promiscuous) to $\Gamma_\wedge$ (cooperative arrest requires all contacts).

\[\langle D_{\triangle\wedge} \;;\; T_\text{network} \;;\; R_{\supseteq} \;;\; P_{\text{pm\_sym}} \;;\; F_\hbar \;;\; K_\text{trap} \;;\; G_\aleph \;;\; \Gamma_\wedge(\text{BRD}) \;;\; \Phi_\text{sub} \;;\; n:m \rangle\]

The condensate gel encoding is identical to the amyloid synthon from PROTEINS.md. This is not a coincidence --- gelation and amyloid formation share the same primitive signature. The framework encodes their equivalence formally.

\textbf{Distance:} $d = 4.00$. Meet conflict on $\Gamma$ only --- the grammar is the single driver of the liquid-gel transition in the shared primitive space. All other primitives are determined by the ordering: $F_\text{gel} > F_\text{liquid}$ ($F_\ell$ is the meet), $K_\text{trap} > K_\text{fast}$ ($K_\text{fast}$ is the meet), $G_\aleph > G_\gimel$ ($G_\gimel$ is the meet). The gel is unambiguously ''above'' the liquid in the primitive lattice.

\textbf{Directed distance asymmetry:} $d(\text{liquid} \to \text{gel}) = 2.80$; $d(\text{gel} \to \text{liquid}) = 2.50$. Counterintuitively, the gel-to-liquid directed distance is slightly shorter --- but the path is blocked. The F-floor prevents traversal of the gel$\rightarrow$liquid direction because the destination has lower $F$ than the source, and source is at $K_\text{trap}$. Directed distance measures the primitive change required; the HotSwap path constraint determines whether it can be executed. Gelation is thermodynamically downhill in primitive space AND kinetically locked --- a doubly efficient trap.

\subsubsection{Class 4 --- Active Gel}

\textbf{Active gel (actin + myosin + ATP):} The archetypal biological programmable matter. Actin filament network ($D_\triangle$) coupled to the ATP hydrolysis temporal cycle ($D_\infty$) gives $D_{\triangle\infty}$. Directed motor forces make this system polarity-directional ($P_\text{directional}$). Filaments assemble/disassemble on minute timescales: $F_\eth$. ATP hydrolysis $\Delta G^\ddagger \sim 80$ kJ/mol: $K_\text{mod}$. Global cytoskeletal flows and cell-scale organization: $G_\aleph$. Sequential grammar: nucleation $\to$ elongation $\to$ motor action ($\Gamma_\to(\text{SEL})$). Near onset of collective motion (active turbulence): $\Phi_c$.

\[\langle D_{\triangle\infty} \;;\; T_\text{network} \;;\; R_\ddagger \;;\; P_\text{directional} \;;\; F_\eth \;;\; K_\text{mod} \;;\; G_\aleph \;;\; \Gamma_\to(\text{SEL}) \;;\; \Phi_c \;;\; n:m \rangle\]

The active gel is the only material in the set with $D_{\triangle\infty}$ --- the temporal cycle is structurally essential. This has an immediate algebraic consequence: the locked state of active gel conflicts on $D$ with every other material. \textbf{There is no shared locked state between active gel and any other PM class} (P-47). Active gel structural integrity requires the temporal cycle, which requires ATP. A static actin-DNA composite is algebraically incoherent.

The Varma probe on active gel gives score 0.60 (''approaching $\Phi_c$'') and degeneracy strength 0.40 --- consistent with the experimental observation that active turbulence onset is a critical phenomenon with power-law velocity correlations.

\subsubsection{Class 5 --- Shape-Memory Polymer}

\textbf{SMP rigid ($T < T_g$):} Crystalline switching domains lock the temporary shape. Crystalline domains function as effectively permanent contacts below $T_g$: $F_\hbar$. Tg-crossing requires $\Delta G^\ddagger > 100$ kJ/mol: $K_\text{slow}$. Mesoscale crystalline domains (10--100 nm) define the structural scale: $G_\gimel$. The shape was programmed in a sequence (heat $\rightarrow$ deform $\rightarrow$ cool $\rightarrow$ fix): $\Gamma_\to(\text{SEL})$.

\[\langle D_\triangle \;;\; T_\text{network} \;;\; R_{\supseteq} \;;\; P_{\text{pm\_pseudo}} \;;\; F_\hbar \;;\; K_\text{slow} \;;\; G_\gimel \;;\; \Gamma_\to(\text{SEL}) \;;\; \Phi_\text{sub} \;;\; n:n \rangle\]

\textbf{SMP elastic ($T > T_g$):} Crystalline domains melted; entropy-elastic recovery drives shape return. Entropic fidelity: $F_\eth$. Thermal energy at $T > T_g$ sufficient to drive recovery: $K_\text{mod}$. All network connections participate simultaneously in elastic recovery: $\Gamma_\wedge(\text{SEL})$.

\[\langle D_\triangle \;;\; T_\text{network} \;;\; R_{\supseteq} \;;\; P_{\text{pm\_pseudo}} \;;\; F_\eth \;;\; K_\text{mod} \;;\; G_\gimel \;;\; \Gamma_\wedge(\text{SEL}) \;;\; \Phi_\text{sub} \;;\; n:n \rangle\]

\textbf{Distance:} $d = 1.70$ --- the closest programmability pair. Meet conflict on $\Gamma$ only: the switch from sequential programming grammar ($\Gamma_\to$, heat-deform-cool-fix) to simultaneous recovery grammar ($\Gamma_\wedge$) is the entire SMP mechanism. Both states share $D_\triangle$, $T_\text{network}$, $G_\gimel$ --- the material's identity (crosslinked polymer with mesoscale crystalline domains) is preserved completely. Only the grammar changes. The SMP mechanism is a grammar transition.

The path search confirms the asymmetry: SMP elastic $\to$ rigid is found in 1 hop ($\Delta\xi_{CP} = +0.611$ nat); SMP rigid $\to$ elastic is blocked. Cooling below $T_g$ is easier to traverse in the primitive graph than heating above it --- reflecting the higher kinetic accessibility of the crystallization direction.

\subsubsection{Class 6 --- Liquid Crystal}

\textbf{LC nematic ($T < T_{NI}$):} Orientational order along the director; no positional order. Molecules aligned: $T_\text{linear}$. Orientational order is moderate (easily disrupted by thermal or field): $F_\eth$. Director reorientation fast under field (sub-ms): $K_\text{fast}$. Mesoscale domain order (correlation length \textasciitilde{}$\mu$m): $G_\gimel$. All neighbours participate in orientation: $\Gamma_\wedge(\text{BROAD})$.

\[\langle D_\triangle \;;\; T_\text{linear} \;;\; R_{\supseteq} \;;\; P_{\text{pm\_sym}} \;;\; F_\eth \;;\; K_\text{fast} \;;\; G_\gimel \;;\; \Gamma_\wedge(\text{BRD}) \;;\; \Phi_\text{sub} \;;\; n:n \rangle\]

\textbf{LC isotropic ($T > T_{NI}$):} Orientational order destroyed. Only local packing correlations: $G_\beth$. Any orientation equally valid: $\Gamma_\vee(\text{BROAD})$. No persistent preference: $F_\ell$.

\[\langle D_\triangle \;;\; T_\text{network} \;;\; R_{\supseteq} \;;\; P_{\text{pm\_sym}} \;;\; F_\ell \;;\; K_\text{fast} \;;\; G_\beth \;;\; \Gamma_\vee(\text{BRD}) \;;\; \Phi_\text{sub} \;;\; n:n \rangle\]

\textbf{Distance:} $d = 3.10$. Meet conflicts on $T$ and $\Gamma$: both the topology (aligned linear vs disordered network) and the grammar (cooperative vs promiscuous) change simultaneously. The $T$-conflict encodes a first-order transition. The LC nematic$\rightarrow$isotropic transition has latent heat (\textasciitilde{}2 kJ/mol) and hysteresis --- consistent with the T-conflict prediction (P-40). The HotSwap path is blocked: ''D/T mismatch --- D\_triangle/T\_linear $\neq$ D\_triangle/T\_network.'' The N$\rightarrow$I transition cannot be traversed in the HotSwap graph; it requires a topology discontinuity.

The directed distance asymmetry: $d(\text{nematic} \to \text{isotropic}) = 3.10$; $d(\text{isotropic} \to \text{nematic}) = 2.50$. Ordering from the isotropic phase (smaller directed distance) is easier than disordering from the nematic --- consistent with the known phenomenon that isotropic-to-nematic alignment occurs faster under an applied field than the reverse.

\begin{center}
\rule{0.5\textwidth}{0.4pt}
\end{center}

\subsection{Part II: The Programmability Algebra}

\subsubsection{Distance Matrix and the Programmability Pairs}

The full $11 \times 11$ distance matrix has the structure expected from the encodings: colloidal fluid and LC isotropic are nearest ($d = 0.00$ --- identical primitive tuples for different materials); active gel is most isolated (maximum distances to all other materials); condensate liquid and gel are at $d = 4.00$; SMP pair at $d = 1.70$.

The programmability pair distances rank the switching difficulty:

\begin{table}[htbp]
\centering
\small
\rowcolors{1}{white}{white}
\begin{tabularx}{\textwidth}{>{\centering\arraybackslash}X >{\centering\arraybackslash}X >{\centering\arraybackslash}X >{\centering\arraybackslash}X >{\centering\arraybackslash}X}
\toprule
\rowcolor{gray!25}\textbf{Pair} & \textbf{$d_\text{sym}$} & \textbf{$d(A \to B)$} & \textbf{$d(B \to A)$} & \textbf{Meet conflicts} \\
\midrule
\rowcolors{1}{white}{gray!10}
SMP rigid $\leftrightarrow$ elastic & 1.70 & 1.20 & 1.10 & $\Gamma$ \\
LC nematic $\leftrightarrow$ isotropic & 3.10 & 3.10 & 2.50 & $T, \Gamma$ \\
Condensate liquid $\leftrightarrow$ gel & 4.00 & 2.80 & 2.50 & $\Gamma$ \\
Colloidal crystal $\leftrightarrow$ fluid & 5.10 & 4.10 & 3.90 & $T, \Gamma$ \\
DNA origami $\leftrightarrow$ strand disp. & 6.10 & 5.60 & 5.50 & $D, T, \Gamma$ \\
\bottomrule
\end{tabularx}
\end{table}

\emph{Prediction (P-38):} This rank predicts the switching energy hierarchy: $\Delta H_\text{SMP} < \Delta H_\text{LC} < \Delta H_\text{colloidal}$.

\subsubsection{Path Searches in the HotSwap Graph}

Eight reconfiguration queries from \texttt{find\_path()}:

\begin{table}[htbp]
\centering
\small
\rowcolors{1}{white}{white}
\begin{tabularx}{\textwidth}{>{\centering\arraybackslash}X >{\centering\arraybackslash}X >{\centering\arraybackslash}X}
\toprule
\rowcolor{gray!25}\textbf{Transition} & \textbf{Result} & \textbf{Notes} \\
\midrule
\rowcolors{1}{white}{gray!10}
SMP rigid $\rightarrow$ elastic ($T_g$ crossing) &  blocked & No intermediate in D\_triangle/T\_network cluster \\
SMP elastic $\rightarrow$ rigid (inverse) &  1 hop, $\Delta\xi = +0.611$ nat & Crystallization is the accessible direction \\
LC nematic $\rightarrow$ isotropic &  blocked & D/T mismatch ($T_\text{lin} \neq T_\text{net}$) \\
LC isotropic $\rightarrow$ nematic &  blocked & D/T mismatch \\
Condensate liquid $\rightarrow$ gel &  1 hop, $\Delta\xi = +0.568$ nat & Gelation is thermodynamically downhill \\
Condensate gel $\rightarrow$ liquid &  blocked & F-floor: $F_\hbar$ at $K_\text{trap}$ \\
Colloidal crystal $\rightarrow$ fluid &  blocked & D/T mismatch ($T_\text{net\_sym} \neq T_\text{net}$) \\
Colloidal fluid $\rightarrow$ crystal &  blocked & D/T mismatch \\
\bottomrule
\end{tabularx}
\end{table}

\textbf{The pattern is striking.} Of eight reconfiguration queries, only two find paths --- both in the ''locking'' direction (elastic$\rightarrow$rigid, liquid$\rightarrow$gel). Every ''unlocking'' transition is blocked. The framework encodes a universal asymmetry: \textbf{locking is downhill in primitive space; unlocking requires external agency that changes $F$ or overcomes $K_\text{trap}$.}

The two systems with $\Phi_c$ (condensate and active gel) are the only ones where the locking direction is also the biologically pathological direction (amyloid formation, gel arrest). The framework encodes the disease-progression direction as the thermodynamically accessible path.

\subsubsection{Primitive Jacobian}

For each pair, which primitive is the dominant design lever? Computed by perturbing $F$, $K$, $G$ one at a time in the rigid/locked state and measuring $\partial d / \partial \text{primitive}$:

\begin{table}[htbp]
\centering
\small
\rowcolors{1}{white}{white}
\begin{tabularx}{\textwidth}{>{\centering\arraybackslash}X >{\centering\arraybackslash}X >{\centering\arraybackslash}X >{\centering\arraybackslash}X >{\centering\arraybackslash}X}
\toprule
\rowcolor{gray!25}\textbf{Pair} & \textbf{Dominant lever} & \textbf{$\Delta d$} & \textbf{Secondary} & \textbf{$\Delta d$} \\
\midrule
\rowcolors{1}{white}{gray!10}
DNA origami $\rightarrow$ strand disp. & $F$ (HIGH$\rightarrow$MEDIUM) & $-$1.10 & $K$ (SLOW$\rightarrow$MODERATE) & $-$1.00 \\
Colloidal crystal $\rightarrow$ fluid & $F$ (HIGH$\rightarrow$LOW) & $-$1.20 & $K$ (SLOW$\rightarrow$FAST) & $-$1.00 \\
Condensate gel $\rightarrow$ liquid & $K$ (TRAP$\rightarrow$FAST) & $-$1.50 & $F$ (HIGH$\rightarrow$LOW) & $-$1.20 \\
SMP rigid $\rightarrow$ elastic & $F$ (HIGH$\rightarrow$MEDIUM) & $-$0.60 & $K$ (SLOW$\rightarrow$MODERATE) & $-$0.50 \\
LC nematic $\rightarrow$ isotropic & $F$ (MEDIUM$\rightarrow$LOW) & $-$0.60 & $G$ (MESOSCALE$\rightarrow$LOCAL) & $-$0.40 \\
\bottomrule
\end{tabularx}
\end{table}

$F$ dominates 4/5 pairs. The exception is condensate gel: $K$ dominates because escaping $K_\text{trap}$ requires more primitive change than lowering $F$ from $F_\hbar$ to $F_\ell$. The practical consequence: for most programmable matter design problems, the fidelity primitive is the first-order design variable. For therapeutic dissolution of biological gels, $K$-targeting (disaggregase) gives the largest primitive leverage.

Prediction P-46: engineering the Jacobian-dominant primitive alone achieves $> 40\%$ of the maximum possible d-reduction.

\subsubsection{Tensor Products --- Composite Programmable Matter}

\begin{table}[htbp]
\centering
\small
\rowcolors{1}{white}{white}
\begin{tabularx}{\textwidth}{>{\centering\arraybackslash}X >{\centering\arraybackslash}X >{\centering\arraybackslash}X}
\toprule
\rowcolor{gray!25}\textbf{Composite} & \textbf{Notation} & \textbf{Key emergent properties} \\
\midrule
\rowcolors{1}{white}{gray!10}
active\_gel $\otimes$ DNA strand disp & $\langle D_\wedge; T_\text{net}; R_{\supseteq}; P_{\text{pm\_pseudo}}; F_\eth; K_\text{mod}; G_{\aleph}; \Gamma_\to(\text{SEL}); \Phi_c \rangle$ & $\Phi_c$ propagates; global coordination with sequence specificity; $G_{\aleph}$ from active gel \\
condensate $\otimes$ SMP elastic & $\langle D_\triangle; T_\text{net}; R_{\supseteq}; P_{\text{pm\_pseudo}}; F_\ell; K_\text{mod}; G_{\gimel}; \Gamma_\wedge(\text{BRD}); \Phi_c \rangle$ & $F_\ell$ bottleneck; $\Phi_c$ propagates; temperature-responsive with critical behavior \\
DNA origami $\otimes$ colloidal crystal & --- & $F_\hbar$ preserved; $G_{\aleph}$ dominant; $T$ conflict between $T_\text{net}$ and $T_{\text{net\_sym}}$ \\
active\_gel $\otimes$ condensate liquid & --- & $\Phi_c$ join-dominant; $G_{\aleph}$ dominant; both $D_{\triangle\infty}$ and $D_{\triangle\wedge}$ --- temporal cycle + mesoscale droplet \\
\bottomrule
\end{tabularx}
\end{table}

The tensor rule $\Phi_c$ join-dominant: any composite containing at least one $\Phi_c$ component inherits $\Phi_c$. This is the algebraic basis for P-41 and P-42: the near-critical component dominates the global responsiveness of the composite.

\subsubsection{The Programmability Lattice}

The meet of all dynamic states (the ''dynamic floor'') and the join of all locked states (the ''rigid ceiling'') bracket the programmable matter design space.

\textbf{Dynamic floor} (meet of six fluid/dynamic systems):

\[\langle D_\triangle \;;\; T_\text{network} \;;\; R_{\supseteq} \;;\; P_{\text{pm\_sym}} \;;\; F_\ell \;;\; K_\text{mod} \;;\; G_\beth \;;\; \Gamma_\vee(\text{BROAD}) \;;\; \Phi_c \rangle\]

The $\Phi_c$ assignment in the floor is emergent --- it propagates from condensate liquid and active gel through the meet lattice (join-dominant). No individual encoding was assigned $\Phi_c$ because criticality was assumed to be optimal; each $\Phi_c$ assignment is physically motivated (near spinodal point, near onset of collective motion). That the meet carries $\Phi_c$ is a theorem, not an assumption.

\textbf{Rigid ceiling} (join of four locked systems):

\[\langle \top \;;\; T_\text{network} \;;\; R_{\supseteq} \;;\; \top \;;\; F_\hbar \;;\; K_\text{slow} \;;\; G_\aleph \;;\; \top \;;\; \Phi_\text{sub} \rangle\]

Conflicts on $D$, $P$, $\Gamma$: the locked states of different materials do not share a single locking mechanism. DNA origami locks by specific Watson-Crick grammar ($\Gamma_\wedge(\text{SPEC})$); colloidal crystal by broad cooperativity ($\Gamma_\wedge(\text{BRD})$); SMP by sequential programming ($\Gamma_\to(\text{SEL})$). No universal lock exists. Any agent targeting ''lock'' in general will be ineffective against all material classes; only agents matched to a specific material's $\Gamma$ will lock that material selectively.

\textbf{Programmability gap:} $d(\text{dynamic floor} \to \text{rigid ceiling}) = 5.40$ nats. The gap is traversed by simultaneously changing $F$, $K$, $G$ --- the three primitives that separate the floor from the ceiling without conflict. These three are the design axes of programmable matter engineering.

\begin{center}
\rule{0.5\textwidth}{0.4pt}
\end{center}

\subsection{Part III: Cross-Domain Analogies}

The catalog metric generates structural identifications between programmable matter states and previously encoded synthons from other domains. These are not hand-selected analogies --- they are the nearest-neighbor outputs of the distance function on the full catalog.

\subsubsection{Condensate liquid $\approx$ Allosteric domain ($d = 2.50$)}

The nearest catalog neighbor to condensate\_liquid is the allosteric\_domain synthon --- closer than any molecular or quantum synthon, closer than any other programmable matter state. The shared primitive cluster: $F_\ell$ (individually weak contacts), $K_\text{fast}$ (dynamic exchange), $G_\gimel$ (mesoscale scale of control), $\Gamma_\vee$ (promiscuous), $\Phi_c$ (near-critical). This is precisely the tuple structure of a mesoscale allosteric system.

The structural identification predicts: condensates implement mesoscale allostery (P-43). A signal entering the condensate (phosphorylation of a client, addition of a competing binder) propagates to the entire droplet via G/D degeneracy --- not just to the immediate vicinity of the signal, but to all clients within the droplet. The effective Hill coefficient of condensate-mediated catalysis is predicted to match the Varma criticality score to within 15\%: $n_H \approx 1.2$.

\subsubsection{Active gel $\approx$ Allosteric domain ($d = 3.40$), Tide pool ecosystem ($d = 4.10$)}

The active gel is structurally nearest to the allosteric\_domain synthon (at greater distance than condensate) but also near the tide\_pool\_ecological synthon --- an ecological-scale analog. The shared structure with the ecological synthon: $D_{\triangle\infty}$ (temporal cycle), $G_\aleph$ (global coordination), $\Phi_c$ (near critical). Active gels implement the same primitive structure as ecological systems that maintain themselves far from equilibrium through continuous energy dissipation. The identification is not a metaphor --- it is a statement about primitive equivalence: the ecological system and the active gel are the same type of constraint-propagation machine at different scales.

\subsubsection{Colloidal crystal $\approx$ Topological insulator ($d = 2.80$)}

The shared structure: $F_\hbar$ (high-fidelity lattice binding), $G_\aleph$ (global crystalline order), $T_\text{network}$ (lattice topology), $R_{\supseteq}$ (non-covalent recognition). The topological insulator and the colloidal crystal are the same primitive system. The framework predicts that colloidal crystals with the appropriate interaction symmetry support topologically protected boundary states (P-44). This prediction is confirmed for DNA-coated colloidal crystals (Rechtsman 2016 and subsequent). The framework extends the boundary: the criterion is $d < 3.0$ from topological\_insulator. All colloidal crystals satisfying this metric with tuned interaction symmetry should exhibit topological boundary states.

\subsubsection{DNA origami $\approx$ Topological insulator ($d = 3.70$)}

More distant than colloidal crystal, but still within the topological analogy range. DNA origami structures with $\Omega \neq 0$ (knots, catenanes, Borromean rings) are predicted to show polynomial rather than exponential strand-failure sensitivity: topological protection replaces the exponential sensitivity of open structures with polynomial scaling (P-45).

\begin{center}
\rule{0.5\textwidth}{0.4pt}
\end{center}

\subsection{Part IV: Predictions Summary}

Ten formal predictions (P-38 through P-47), all algebraically confirmed, all requiring experimental validation:

\begin{table}[htbp]
\centering
\small
\rowcolors{1}{white}{white}
\begin{tabularx}{\textwidth}{>{\centering\arraybackslash}X >{\centering\arraybackslash}X >{\centering\arraybackslash}X >{\centering\arraybackslash}X}
\toprule
\rowcolor{gray!25}\textbf{ID} & \textbf{Domain} & \textbf{Prediction} & \textbf{Status} \\
\midrule
\rowcolors{1}{white}{gray!10}
P-38 & SMP vs LC vs colloidal & $\Delta H_\text{SMP} < \Delta H_\text{LC} < \Delta H_\text{colloidal}$ from $d$ rank & pending calorimetry comparison \\
P-39 & Condensate gel rescue & Thermal stimulus alone insufficient; requires $\Gamma$- or K-targeting &  path blocked; consistent with Hsp70 biology \\
P-40 & LC and colloidal transitions & T-conflict $\rightarrow$ first-order (latent heat, hysteresis) &  both known to be first-order \\
P-41 & Actin-DNA composite & Collective motion only with ATP; $\Phi_c$ ATP-conditional & pending not tested \\
P-42 & Programmable matter universality & Maximal versatility $\rightarrow$ near-critical (dynamic floor = $\Phi_c$) &  consistent with cortex/cytoplasm criticality \\
P-43 & Condensate allostery & Condensates implement mesoscale allostery; $n_H \approx 1.2$ & pending Hill coefficient prediction novel \\
P-44 & Topological colloids & $d < 3.0$ from TI $\rightarrow$ boundary states with tuned symmetry &  Rechtsman 2016; range extended \\
P-45 & Topological DNA origami & Knots/catenanes show polynomial (not exponential) error scaling & pending not yet tested \\
P-46 & PM design Jacobian & Dominant Jacobian primitive achieves $> 40\%$ of full optimization & pending Tg vs crosslink experiment pending \\
P-47 & Actin-DNA co-programmability & Actin-DNA cannot be jointly locked; structural integrity requires ATP & pending consistent with actin biology \\
\bottomrule
\end{tabularx}
\end{table}

\begin{center}
\rule{0.5\textwidth}{0.4pt}
\end{center}

\subsection{Part V: What the Framework Cannot Say About Programmable Matter}

The same epistemic discipline applied in METAPHYSICS.md applies here.

\textbf{The framework does not predict switching speed.} $K$ encodes the ordinal kinetic tier (fast/moderate/slow/trap), not the absolute rate. Whether the switching kinetics of a specific SMP are 10 ms or 10 s requires domain equations (Williams-Landel-Ferry, viscosity models) that the framework deliberately excludes.

\textbf{The framework does not predict mechanical properties.} Stiffness, Young's modulus, and toughness are continuous scalars. The framework's primitives encode the type of structural order ($T$, $G$) and the fidelity of contacts ($F$), not the quantitative mechanical response.

\textbf{The framework does not predict processing conditions.} What temperature, concentration, or field strength is required to achieve a given state change is outside the grammar's scope. The grammar says \emph{whether} a state change is accessible; it says nothing about the conditions under which it occurs.

\textbf{The framework does not guarantee that tensor composites are synthetically accessible.} The algebra computes the primitive tuple of a composite; it does not guarantee that the composite can be made. $\text{active\_gel} \otimes \text{DNA strand disp.}$ yields a prediction about what properties an actin-DNA hybrid would have --- it does not indicate whether such a hybrid is stable, scalable, or useful.

What the framework does do, in this domain as in others: it establishes the relational structure of programmable matter. Programmability is not a property of a material; it is a property of a pair of states connected by a path in primitive space. The design question is a graph-theoretic question. The grammar makes that question precise.

\begin{center}
\rule{0.5\textwidth}{0.4pt}
\end{center}

\emph{See \texttt{programmable\_matter\_tests1.py} and \texttt{programmable\_matter\_tests2.py} for the exact algebra computations underlying all results in this document. Formal predictions are registered in \texttt{PRIMITIVE\_PREDICTIONS.md} (P-38 through P-47). The theoretical framework is in \texttt{SYNTHONICON.md} §X.}


\end{document}
